% Appendix
% Paper: "Just a Glimpse: Temporal Wearing as a Design Paradigm for Transient Mixed Reality"

\appendix

\section{Prototype Descriptions}
\label{appendix:prototypes}

This appendix provides detailed descriptions of all six temporal wearing prototypes. Each description includes the design rationale, technical specifications, form factor characteristics, and transition profile.

\subsection{Prototype A: [PLACEHOLDER Name] (Handheld)}

\begin{figure}[h]
    \centering
    \includegraphics[width=0.8\columnwidth]{figures/prototype-a-placeholder}
    \caption{Prototype A: [PLACEHOLDER]. A handheld monocular/binocular viewer designed for single-hand operation. [PLACEHOLDER: brief description of appearance and use gesture.]}
    \label{fig:prototype-a}
\end{figure}

\paragraph{Design Rationale.}
[PLACEHOLDER: 2--3 sentences on why this form factor was chosen, what design space position it occupies, and what historical precedent it draws from (e.g., lorgnette, magnifying glass).]

\paragraph{Technical Specifications.}
\begin{itemize}
    \item Display: [PLACEHOLDER: type, resolution, FOV]
    \item Tracking: [PLACEHOLDER: method]
    \item Weight: [PLACEHOLDER] g
    \item Dimensions: [PLACEHOLDER] mm
    \item Power: [PLACEHOLDER: battery type, life]
    \item Onset time: [PLACEHOLDER] s (measured)
    \item Offset time: [PLACEHOLDER] s (measured)
\end{itemize}

\paragraph{Transition Profile.}
[PLACEHOLDER: Description of the don/doff gesture, what makes it fast, any design features that support instant-on.]

\subsection{Prototype B: [PLACEHOLDER Name] (Handheld)}

\begin{figure}[h]
    \centering
    \includegraphics[width=0.8\columnwidth]{figures/prototype-b-placeholder}
    \caption{Prototype B: [PLACEHOLDER]. [PLACEHOLDER: brief description.]}
    \label{fig:prototype-b}
\end{figure}

[PLACEHOLDER: Same structure as Prototype A --- design rationale, technical specifications, transition profile.]

\subsection{Prototype C: [PLACEHOLDER Name] (Neck-hung)}

\begin{figure}[h]
    \centering
    \includegraphics[width=0.8\columnwidth]{figures/prototype-c-placeholder}
    \caption{Prototype C: [PLACEHOLDER]. A neck-hung viewer suspended from a lanyard, resting on the chest when not in active use. [PLACEHOLDER: brief description.]}
    \label{fig:prototype-c}
\end{figure}

\paragraph{Design Rationale.}
[PLACEHOLDER: Rationale for neck-hung form factor. How passive body coupling enables rapid transition from rest to active use. Historical precedent: press photographer cameras, tourist binoculars worn on neck straps.]

\paragraph{Technical Specifications.}
\begin{itemize}
    \item Display: [PLACEHOLDER]
    \item Tracking: [PLACEHOLDER]
    \item Weight: [PLACEHOLDER] g
    \item Lanyard: [PLACEHOLDER: material, length, adjustment]
    \item Onset time: [PLACEHOLDER] s
    \item Offset time: [PLACEHOLDER] s
\end{itemize}

\paragraph{Transition Profile.}
[PLACEHOLDER: Description of the raise-from-chest gesture, how gravity assists offset.]

\subsection{Prototype D: [PLACEHOLDER Name] (Neck-hung)}

\begin{figure}[h]
    \centering
    \includegraphics[width=0.8\columnwidth]{figures/prototype-d-placeholder}
    \caption{Prototype D: [PLACEHOLDER]. [PLACEHOLDER: brief description.]}
    \label{fig:prototype-d}
\end{figure}

[PLACEHOLDER: Same structure --- design rationale, technical specifications, transition profile.]

\subsection{Prototype E: [PLACEHOLDER Name] (Stationary)}

\begin{figure}[h]
    \centering
    \includegraphics[width=0.8\columnwidth]{figures/prototype-e-placeholder}
    \caption{Prototype E: [PLACEHOLDER]. A plinth-mounted or wall-mounted viewfinder that users lean into. [PLACEHOLDER: brief description.]}
    \label{fig:prototype-e}
\end{figure}

\paragraph{Design Rationale.}
[PLACEHOLDER: Rationale for stationary form factor. The body approaches the device rather than vice versa. Historical precedent: coin-operated viewfinders, museum stereoscopes, peep shows. How fixed placement maximizes social legibility and eliminates device management.]

\paragraph{Technical Specifications.}
\begin{itemize}
    \item Display: [PLACEHOLDER]
    \item Tracking: [PLACEHOLDER]
    \item Mounting: [PLACEHOLDER: plinth/wall/pedestal, height, adjustability]
    \item Viewing aperture: [PLACEHOLDER: dimensions, eye relief]
    \item Onset time: [PLACEHOLDER] s
    \item Offset time: [PLACEHOLDER] s
\end{itemize}

\paragraph{Transition Profile.}
[PLACEHOLDER: Description of the lean-in gesture, proximity-triggered activation, natural withdrawal.]

\subsection{Prototype F: [PLACEHOLDER Name] (Stationary)}

\begin{figure}[h]
    \centering
    \includegraphics[width=0.8\columnwidth]{figures/prototype-f-placeholder}
    \caption{Prototype F: [PLACEHOLDER]. [PLACEHOLDER: brief description.]}
    \label{fig:prototype-f}
\end{figure}

[PLACEHOLDER: Same structure --- design rationale, technical specifications, transition profile.]


\section{Interview Protocol}
\label{appendix:interview}

Semi-structured exit interview guide used with visitors following interaction with temporal wearing prototypes. Interviews lasted 10--15 minutes and were audio-recorded with consent.

\subsection{Opening}
\begin{enumerate}
    \item Thank you for agreeing to speak with us. Can you tell me a bit about your visit today? Who are you here with?
    \item I noticed you used [prototype name/description]. Can you walk me through what happened, from when you first noticed it?
\end{enumerate}

\subsection{Discovery and Approach}
\begin{enumerate}[resume]
    \item What first caught your attention about the device?
    \item What did you think it was before you used it?
    \item Was there anything that made you hesitate before picking it up / leaning in?
\end{enumerate}

\subsection{The Experience}
\begin{enumerate}[resume]
    \item Can you describe what you saw when you looked through it?
    \item How long do you think you looked through it the first time?
    \item Did you look through it more than once? [If yes:] What made you want to look again?
    \item How did it feel to put it down / step away?
\end{enumerate}

\subsection{Social Dimensions}
\begin{enumerate}[resume]
    \item Were you aware of other people around you while you were using it?
    \item Did you feel self-conscious at all? [Probe: why/why not?]
    \item Did you share the experience with anyone? [If yes:] How?
    \item [If in group:] Did your companion(s) also try it? How did that come about?
\end{enumerate}

\subsection{Comparison and Reflection}
\begin{enumerate}[resume]
    \item Have you used VR or AR devices before? [If yes:] How did this compare?
    \item If this device were available at other places you visit---museums, parks, stores---would you use it?
    \item Is there anything else you'd like to share about the experience?
\end{enumerate}


\section{Survey Instruments}
\label{appendix:survey}

\subsection{User Social Comfort Survey}

Administered immediately after interaction. All items rated on a 7-point Likert scale (1 = Strongly Disagree, 7 = Strongly Agree).

\begin{enumerate}
    \item I felt comfortable using this device in this setting.
    \item I was aware of other people watching me while I used the device.
    \item I felt self-conscious while using the device. [R]
    \item The act of picking up / leaning into the device felt natural.
    \item I would feel comfortable using this device in a different public setting (e.g., a park, a shop).
    \item I would use this device again if I encountered it.
    \item The device was easy to start using (no instructions needed).
    \item I felt socially disconnected from my companions while using the device. [R]
\end{enumerate}

Items marked [R] are reverse-coded.

\subsection{Bystander Comfort Survey}

Administered to companions or nearby visitors who observed but did not use the device.

\begin{enumerate}
    \item I could tell what the person was doing with the device.
    \item The person's use of the device seemed natural in this setting.
    \item I felt uncomfortable watching the person use the device. [R]
    \item I understood that the interaction would be brief.
    \item I felt the person was still ``present'' and available during their use of the device.
    \item Seeing the person use the device made me want to try it myself.
\end{enumerate}


\section{Video Coding Scheme}
\label{appendix:coding}

Adapted from Brignull and Rogers~\cite{brignull2003enticing} and extended for temporal wearing contexts. Each interaction episode (defined as a continuous sequence from first approach to final departure from the device vicinity) was coded on the following dimensions.

\subsection{Discovery Phase}
\begin{itemize}
    \item \textbf{Discovery trigger}: Self-initiated (noticed device independently) / Social (prompted by companion or observed another user) / Signage (read label or sign) / Staff (directed by venue staff)
    \item \textbf{Approach trajectory}: Direct (walked straight to device) / Circuitous (approached indirectly, pausing at other exhibits) / Return (came back after initial pass)
    \item \textbf{Hesitation}: None / Brief pause ($<$3s) / Extended pause ($>$3s) / Abandoned approach
\end{itemize}

\subsection{Donning Phase}
\begin{itemize}
    \item \textbf{Donning gesture}: One-hand lift / Two-hand lift / Lean-in / Raise from chest / Other
    \item \textbf{Onset confidence}: Immediate (no adjustment) / Minor adjustment (1 repositioning) / Multiple adjustments / Required assistance
    \item \textbf{Onset time}: Measured in seconds from first contact to stable viewing position
\end{itemize}

\subsection{Glimpse Phase}
\begin{itemize}
    \item \textbf{Glimpse duration}: Measured in seconds
    \item \textbf{Viewing behavior}: Stationary gaze / Scanning (moving device or head) / Searching (active exploration) / Showing companion
    \item \textbf{Verbal behavior during glimpse}: Silent / Exclamation / Narration to companion / Question to companion
    \item \textbf{Body posture}: Upright / Leaning / Crouching / Other
\end{itemize}

\subsection{Doffing Phase}
\begin{itemize}
    \item \textbf{Doffing gesture}: Lower to surface / Lower to chest / Hand to companion / Step back / Other
    \item \textbf{Offset time}: Measured in seconds
    \item \textbf{Facial expression at doff}: Neutral / Smile / Surprise / Puzzlement / Other
\end{itemize}

\subsection{Post-Doff Phase}
\begin{itemize}
    \item \textbf{Immediate action}: Verbal sharing / Gestural sharing (pointing) / Re-engagement (picks up again) / Moves on / Waits for companion
    \item \textbf{Social interaction}: None / Narration / Discussion / Hand-off / Joint re-engagement
    \item \textbf{Return}: No return / Return within 1 minute / Return within 5 minutes / Return after 5+ minutes
\end{itemize}

\subsection{Temporal Wearing Behaviors}
\begin{itemize}
    \item \textbf{Double-take}: Quick initial glimpse ($<$3s) followed by re-engagement within 10s
    \item \textbf{Hand-off}: Device passed directly to companion within 15s of doffing
    \item \textbf{Narrator}: User verbally describes augmented content to companion during glimpse
    \item \textbf{Orbit}: User returns to device after visiting other parts of venue ($>$1 minute interval)
\end{itemize}
